\documentclass[11pt]{aime}
\usepackage{lmodern}

% ─────────────────────────────────────────────────────────────────────
% Metadata
% ─────────────────────────────────────────────────────────────────────

\author[1]{Richard Beckmann}
\author[1]{Hannes Öhman}
\author[1]{Pablo Martínez Crespo}
\author[3]{Robert S. Jordan}
\author[4]{Marisa Gliege}
\author[5]{Santiago Miret}
\author[6]{Vijay Kris Narasimhan}
\author[1,2]{Rocío Mercado}
\affiliation[1]{AI Laboratory for Molecular Engineering (AIME), Department of Computer Science and Engineering, Chalmers University of Technology \& University of Gothenburg, Gothenburg, Sweden}
\affiliation[2]{Science for Life Laboratory (SciLifeLab), Gothenburg, Sweden}
\affiliation[3]{Technology Research, Intel Corporation}
\affiliation[4]{Chief Technology Office, EMD Electronics}
\affiliation[5]{Lila Sciences}
\affiliation[6]{Merck KGaA, Darmstadt, Germany}

% Working title -- revisit once Results/Discussion are drafted.
\title{Uncertainty- and Plausibility-Aware Reinforcement Learning for De Novo Surfactant Design}

\contribution[*]{Corresponding author}  % Generally Rocío

\correspondence{\email{rocio.mercado@chalmers.se}}
\codeavailability{\url{https://github.com/hannesohman/REINVENT4Surfactants}}
\date{\today}

\abstract{%
% TODO: draft after Results/Discussion are written.
}

\begin{document}
\maketitle

% =====================================================================
% INTRODUCTION
% =====================================================================
\section{Introduction}
\label{sec:introduction}

Surfactants -- amphiphilic molecules that combine a polar, water-soluble head group with a nonpolar, oil-soluble tail -- are among the highest-volume and most functionally diverse classes of specialty chemicals in use today, underpinning household and industrial detergents, personal-care and cosmetic formulations, emulsification and enhanced oil recovery in the petrochemical industry \cite{jia_review_2024}, agrochemical adjuvants \cite{kovalchuk_surfactant-mediated_2021}, and pharmaceutical excipients \cite{bhadani_current_2020,rosen_surfactants_2012}. Their function follows directly from their chemical character: dissolved surfactant monomers preferentially adsorb at interfaces (air--water, oil--water, solid--liquid), lowering interfacial tension, and above a compound-specific critical micelle concentration (CMC) self-assemble into micelles and higher-order aggregates \cite{rosen_surfactants_2012}. The two properties central to this behaviour -- the CMC (commonly reported as pCMC $=-\log_{10}(\mathrm{CMC})$, where a higher value denotes a more efficient surfactant) and the surface or interfacial tension achieved at or above the CMC -- together largely determine how well a given molecule performs as a wetting agent, emulsifier, foaming agent, or solubiliser, and are the two properties this work targets for optimisation.

Despite this well-established link between molecular structure and interfacial function, rationally designing new surfactants with improved pCMC and surface tension remains difficult, for three compounding reasons. First, environmental and regulatory pressure is actively narrowing the space of acceptable chemistries: certain fluorosurfactant classes are being phased out and require industrial replacements that do not sacrifice performance \cite{sharma_safer_2023}, and more broadly, growing evidence of the persistence and poor biodegradability of conventional, petroleum-derived surfactants is pushing formulators toward alternatives that meet substantially stricter sustainability criteria \cite{bhadani_current_2020,nagtode_green_2023}. Second, the space of chemically plausible amphiphiles -- combinatorially large across head-group chemistry, tail length and branching, and linker chemistry -- vastly exceeds what has been experimentally characterised: measuring CMC and surface tension is comparatively slow and labour-intensive, so curated experimental datasets remain small (on the order of $10^3$ molecules) relative to those routinely available for bioactivity in drug discovery \cite{hodl_surfpro_2025}. Third, the quantitative structure--property relationship (QSPR) models built on these datasets, while increasingly accurate at predicting CMC and related properties for molecules resembling their training data \cite{hu_review_2010,boukelkal_qspr_2024}, are surrogates rather than ground truth: even state-of-the-art machine-learning predictors struggle to generalise to surfactant chemistries or phase behaviour outside their training distribution \cite{thacker_can_2023}, and typically give no signal for when they are being asked to extrapolate.

Machine learning has, over the past decade, begun to reshape how new molecules are proposed to address exactly this kind of problem. Rather than screening pre-existing libraries, generative models learn a distribution over synthesisable chemical space and are steered -- via reinforcement learning (RL), latent-space optimisation, or related strategies -- toward molecules predicted to satisfy one or more target properties \cite{sanchez-lengeling_inverse_2018,ozcelik_generative_2025}. Recurrent-neural-network agents trained with policy-gradient RL, popularised by REINVENT \cite{olivecrona_molecular_2017} and its successors \cite{blaschke_memory-assisted_2020,loeffler_reinvent_2024}, are among the most established instances of this paradigm and have been applied extensively to bioactive small-molecule design. Their extension to functional, non-bioactive molecules -- materials, catalysts, and industrial specialty chemicals -- has so far been comparatively limited. For surfactants specifically, only a single prior study has coupled a generative model with reinforcement learning guided by predicted properties, using a SELFIES-based variational autoencoder paired with a graph-neural-network property predictor \cite{nnadili_surfactant-specific_2024}; that framework, however, optimises a single target property, the critical micelle concentration, in isolation. Genuinely multi-objective surfactant design -- jointly targeting several, potentially competing performance properties, as most real formulation problems require -- together with the RNN/transformer-based REINVENT lineage and the specific problem of designing a reward robust enough for multi-objective RL against learned surrogates, remains largely unexplored for this compound class, despite the same combinatorial-space argument applying at least as strongly as it does for bioactive molecules.

Coupling a generative model to a learned property predictor as its RL reward, however, introduces a specific failure mode: because the reward is only ever as reliable as the surrogate that computes it, an agent can learn to exploit regions where the surrogate is confidently wrong rather than regions of genuinely good chemistry -- a form of reward hacking documented across generative molecular design more broadly \cite{yoshizawa_datadriven_2025}, and one we also observe directly in this system (Section~\ref{sec:methods}). Guarding against this failure mode, while still balancing multiple, sometimes competing, target properties -- here pCMC and surface tension, together with basic requirements such as structural plausibility -- is therefore itself an open methodological question, distinct from generative modelling or property prediction in isolation, and one for which uncertainty-aware training strategies have recently been proposed as a partial answer \cite{medina_uncertainty_2026}.

In this work, we apply REINVENT4 \cite{loeffler_reinvent_2024} to the de novo design of surfactants. A pretrained chemical language model is fine-tuned on experimentally measured surfactant data \cite{hodl_surfpro_2025}, and reinforcement learning is then guided by XGBoost \cite{chen_xgboost_2016} surrogate models of pCMC and surface tension trained on SurfKit, an extended version of this dataset augmenting the curated experimental measurements with molecular-dynamics-derived properties \cite{beckmann_surfkit_2026}. Rather than treating these two properties as a fixed, standalone objective, we explore several complementary strategies for turning them into a robust multi-objective reward: penalising candidates on which an ensemble of property models itself disagrees, rewarding structural resemblance to real, purchasable chemistry independent of the property surrogates, and explicitly tracking the empirical Pareto front of generated molecules over the course of training. We validate the resulting pipeline against two independent sets of real molecules held out from training -- one drawn from experimentally characterised surfactants, one from the purchasable ZINC20 chemical space \cite{irwin_zinc20_2020} -- to confirm that gains in surrogate-predicted score reflect genuine, non-exploited chemistry rather than surrogate blind spots.

Concretely, this paper makes three contributions: (i) we apply REINVENT4 to surfactant molecules, to our knowledge the first application of this RNN/transformer-based generative framework to this class of industrially important, non-bioactive compounds, complementing the SELFIES-VAE-based, single-property (CMC-only) approach explored previously \cite{nnadili_surfactant-specific_2024}; (ii) in contrast to that single-property precedent, we treat surfactant design as a genuinely multi-objective problem, jointly optimising pCMC and surface tension, and we systematically explore and compare several novel strategies for making this multi-objective reward robust to an uncertain, potentially exploitable surrogate -- uncertainty-aware score and loss modulation, structural-plausibility scoring, and Pareto-front tracking -- across a combinatorial sweep of scoring-function configurations; and (iii) building on the best-performing configuration(s), we propose a set of $N$ novel surfactant candidates that our validated surrogates predict to exceed known surfactants in pCMC and/or surface tension.

% =====================================================================
% METHODS
% =====================================================================
\section{Methods}
\label{sec:methods}

\subsection{Generative framework}
\label{sec:methods:generative}

\paragraph{Prior model and transfer learning.}
Molecule generation is performed with REINVENT4\cite{loeffler_reinvent_2024}, using its recurrent neural network prior pretrained on a large subset of PubChem\cite{kim_pubchem_2023}.
Before reinforcement learning (RL), the prior is fine-tuned via transfer learning (TL) on real surfactant SMILES from the SurfKit dataset\cite{beckmann_surfkit_2026} (Section~\ref{sec:methods:holdout}), using SMILES randomisation\cite{bjerrum_smiles_2017} and canonicalisation, per-epoch shuffling, and a decaying learning-rate schedule ($\mathrm{lr}_0=10^{-4}$, decaying to a minimum of $10^{-7}$).
Critically, TL is performed only on the trainval partition of the dataset, excluding the held-out molecules described in Section~\ref{sec:methods:holdout}, so that any subsequent exact rediscovery of a held-out molecule reflects genuine generalisation rather than memorisation during fine-tuning.

\paragraph{Reinforcement learning.}
RL proceeds via REINVENT4's staged-learning mode with the Diversity Augmented Policy (DAP) reward strategy.
At each step, a batch of SMILES is sampled from the agent, scored by the multi-parameter objective described in Sections~\ref{sec:methods:scoring}--\ref{sec:methods:uncertainty}, and used to update the agent via
\begin{equation}
    \mathcal{L}_j = \left(\mathrm{NLL}_\mathrm{prior}(x_j) + \sigma \cdot s(x_j) - \mathrm{NLL}_\mathrm{agent}(x_j)\right)^2,
    \label{eq:dap-loss}
\end{equation}
where $s(x_j)$ is the composite score for molecule $x_j$, $\sigma$ is a scaling hyperparameter controlling how strongly the score influences the augmented likelihood, and $\mathrm{NLL}_\mathrm{prior}$/$\mathrm{NLL}_\mathrm{agent}$ are the negative log-likelihoods under the (fixed) TL checkpoint and the (updating) agent, respectively.
The per-sample losses are averaged and backpropagated via Adam\cite{kingma_adam_2015}.
$\sigma$, the learning rate, and the batch size are treated as hyperparameters and optimised per experimental condition (Section~\ref{sec:methods:hpo}).

\subsection{Surrogate scoring functions}
\label{sec:methods:scoring}

\paragraph{Point-estimate property models.}
Two physicochemical properties form the core of the scoring objective: the critical micelle concentration, expressed as pCMC $= -\log_{10}(\mathrm{CMC})$, and the MD-derived air--water interfacial surface tension (SurfTen).
Point estimates for both are obtained from single XGBoost regressors\cite{chen_xgboost_2016} trained on 50 RDKit\cite{landrum2006rdkit} molecular descriptors, taken from the SurfKit surrogate model suite\cite{beckmann_surfkit_2026} (Section~\ref{sec:methods:holdout}).

\paragraph{Property score transformation and direction convention.}
Because pCMC is defined as $-\log_{10}(\mathrm{CMC})$, a \emph{higher} pCMC corresponds to a \emph{lower} CMC, i.e. a more efficient surfactant that forms micelles at lower concentration.
This was confirmed empirically against the training data: dodecyl sulfate (sodium dodecyl sulfate, SDS) has pCMC $=2.03$ in SurfKit, implying $\mathrm{CMC}=10^{-2.03}\approx9.3$\,mM, matching SDS's well-established experimental CMC.
pCMC is therefore maximised, while SurfTen is minimised.
Each raw prediction $\hat{y}$ is normalised against calibration bounds $[y_\mathrm{min}, y_\mathrm{max}]$ (the 5th/95th percentile of the property over the SurfKit training distribution),
\begin{equation}
    s(\hat{y}) =
    \begin{cases}
        \dfrac{\hat{y}-y_\mathrm{min}}{y_\mathrm{max}-y_\mathrm{min}}, & \text{property maximised (pCMC)}\\[1.2ex]
        1-\dfrac{\hat{y}-y_\mathrm{min}}{y_\mathrm{max}-y_\mathrm{min}}, & \text{property minimised (SurfTen)}
    \end{cases}
    \label{eq:normalize}
\end{equation}
clipped to $[0,1]$, so that a higher transformed score always denotes a more desirable value regardless of the underlying property's own optimisation direction.
Individual component scores are combined via a weighted geometric mean (REINVENT4's multi-parameter optimisation, MPO, scoring type), with weights specified per experimental condition (Section~\ref{sec:methods:design}).

\paragraph{Structural plausibility.}
An objective built only from pCMC and SurfTen point estimates was found to reward molecules the surrogate models score confidently but that are not plausible surfactants (e.g.\ small organosilicon or perfluorinated fragments), since ensemble disagreement (Section~\ref{sec:methods:uncertainty}) reflects inter-fold variance rather than distance from the training distribution.
To discourage this, a structural plausibility term (ZincPlausibility) scores each candidate by the fraction of its BRICS\cite{degen_art_2008} fragments that appear in a reference vocabulary built once from a 200{,}000-molecule uniform random sample of the ZINC20 in-stock (purchasable) chemical space\cite{irwin_zinc20_2020}.
This deliberately avoids molecule-level nearest-neighbour similarity search against the reference set, which would require $O(\text{batch}\times\text{reference})$ Tanimoto comparisons per RL step; fragment-vocabulary lookup instead costs one BRICS decomposition and a set of dictionary lookups per molecule, independent of reference-set size.

\paragraph{Pareto-based scoring.}
Two additional scoring components track the empirical (pCMC, SurfTen) Pareto front\cite{deb_fast_2002} realised by the agent's own previously generated molecules, added alongside (not replacing) the direct pCMC/SurfTen terms.
ParetoBoost assigns a fixed bonus to candidates that dominate or extend the current front;
ParetoGradient instead assigns a continuous score based on the candidate's distance to the nearest segment of the front, so that near-miss candidates are still rewarded proportionally.
Both maintain the front as an incrementally-updated running set (a dominated point can never re-enter the front once superseded, so each step's new candidates only need to be checked against, and folded into, the current front rather than the full generated-molecule history), so the reference front evolves together with the policy at a cost independent of how many steps have already run.

\subsection{Uncertainty-aware scoring: Score Modulation and Loss Modulation}
\label{sec:methods:uncertainty}

Point-estimate XGBoost regressors carry no native uncertainty estimate.
Following Medina and Janet\cite{medina_uncertainty_2026}, uncertainty for pCMC and SurfTen is instead obtained from a 25-member XGBoost ensemble (5 outer cross-validation splits $\times$ 5 fold models per property) trained on the same SurfKit data;
for a given molecule, the standard deviation across the 25 ensemble predictions serves as the uncertainty estimate, following the deep-ensembles approach to predictive uncertainty estimation\cite{lakshminarayanan_simple_2017}.

We implement and compare the two uncertainty-incorporation strategies proposed by Medina and Janet\cite{medina_uncertainty_2026}.
\textbf{Score Modulation (SM)} treats uncertainty as an additional, independent term in the same weighted geometric mean as the property scores,
\begin{equation}
    S_\mathrm{SM}(x) = \mathrm{MPO}\!\left(s_1(x), \ldots, s_K(x), s_\mathrm{unc}(x)\right),
    \label{eq:sm}
\end{equation}
where $s_\mathrm{unc}$ is the (normalised, inverted) ensemble standard deviation, so that low disagreement (reliable prediction) scores near 1 and high disagreement scores near 0.
\textbf{Loss Modulation (LM)} instead leaves the score untouched and reweights each sample's contribution to the RL policy-gradient loss (Eq.~\ref{eq:dap-loss}),
\begin{equation}
    \mathcal{L}_\mathrm{LM} = \frac{1}{N}\sum_{j=1}^{N} \frac{w_j}{\overline{w}}\,\mathcal{L}_j,
    \qquad w_j = \frac{1}{2}\left(s_\mathrm{unc}^{\,\mathrm{pCMC}}(x_j) + s_\mathrm{unc}^{\,\mathrm{SurfTen}}(x_j)\right),
    \label{eq:lm}
\end{equation}
where $\overline{w}$ is the batch mean of $w_j$, so that molecules with lower uncertainty contribute more strongly to the gradient update while the reward signal itself is unaffected.
Because Eq.~\ref{eq:lm} operates on REINVENT4's training loop rather than the scoring function, LM is implemented as a runtime patch of the reward-computation and policy-update routines, applied only when explicitly enabled; with it disabled, REINVENT4's behaviour is unchanged.
SM and LM can be enabled independently or jointly from the same underlying scoring configuration (Section~\ref{sec:methods:design}).

\subsection{Holdout construction and transfer-learning exclusion}
\label{sec:methods:holdout}

\paragraph{SurfKit holdout.}
From the 1551 curated SurfKit molecules, a composite quality score is computed per molecule as
\begin{equation}
    c(x) = \sqrt{s_\mathrm{pCMC}(x)\cdot s_\mathrm{SurfTen}(x)},
    \label{eq:composite}
\end{equation}
using the true measured property values and the same normalisation convention as Eq.~\ref{eq:normalize}.
A flat split that holds out the top-100 molecules by composite score directly was found to leak into transfer learning: SurfKit contains homologous series (the same core scaffold at different alkyl chain lengths), and a flat top-$N$ split has no scaffold-disjointness constraint, so a held-out molecule's exact scaffold could still appear in the trainval set at a different chain length. Sampling directly from a transfer-learning checkpoint trained on such a split, with no reinforcement learning at all, exactly regenerated 65\% of a flat top-100 holdout -- higher than any subsequently RL-optimised combination achieved -- confirming the model was exploiting this leakage rather than generalising.

The holdout is therefore constructed to keep every homologous/structural family entirely on one side of the split. All 1551 molecules are clustered via union-find, merging two molecules if they share an identical non-empty Bemis--Murcko scaffold\cite{bemis_properties_1996} \emph{or} their Morgan-fingerprint\cite{rogers_extended-connectivity_2010} (radius 2, 2048 bit) Tanimoto distance is below 0.35 (fingerprint similarity alone under-merges same-scaffold pairs with a large chain-length difference, since that shifts the fingerprint distance without changing the scaffold). Whole clusters are then greedily added to the holdout, highest mean composite score first, until a target of 100--120 molecules is reached; this yielded 105 holdout molecules from the top 29 (of 161) clusters, with the remaining 1446 molecules forming the transfer-learning trainval set (Section~\ref{sec:methods:generative}). Re-running both diagnostics against the corrected split confirmed the fix: scaffold overlap between holdout and trainval dropped from 81\% to 0\%, and zero-RL-step sampling rediscovery dropped from 65\% to 1.9\%.

\paragraph{ZINC holdout.}
A second, independent holdout is drawn from the ZINC20 in-stock (purchasable) chemical space\cite{irwin_zinc20_2020}: approximately 11.3 million standardised molecules obtained via bulk tranche download, restricted to purchasability classes A/B/C.
Molecules containing any element absent from the SurfKit training vocabulary (C, Cl, F, N, O, P, S, Si) are excluded.
The remaining pool is scored with the same surrogate ensemble and ranked by the identical composite score (Eq.~\ref{eq:composite}); the top 100 molecules form the ZINC holdout.
Because this holdout is ranked purely by surrogate-predicted properties rather than ground truth, it provides a complementary but softer signal than the SurfKit holdout (Section~\ref{sec:methods:evaluation}).

\subsection{Hyperparameter optimisation}
\label{sec:methods:hpo}

Hyperparameters ($\sigma$, learning rate, batch size) are optimised independently for each experimental condition (Section~\ref{sec:methods:design}) using Optuna\cite{optuna_2019} with its default tree-structured Parzen estimator sampler\cite{bergstra_algorithms_2011}, maximising the mean composite \texttt{Score} (Eq.~\ref{eq:sm}/the corresponding MPO for that condition) over the top-5\% scoring molecules of a trial.
The search space is $\sigma\in[50,500]$ (log-uniform, integer), learning rate $\in[10^{-5},10^{-3}]$ (log-uniform), and batch size $\in\{10,50,100,200,500\}$ (categorical); step count is not fixed but derived per trial as a fixed oracle budget of 10{,}000 proposed molecules divided by that trial's batch size, so every trial proposes the same total number of molecules regardless of how it is split between batch size and step count.
An earlier version of this objective took a \emph{fixed} top-100 count rather than a percentile: since the candidate pool size scales with batch size for a fixed step count, a fixed-count top-$N$ objective is an increasingly extreme order statistic for larger batches, mechanically favouring the largest batch size in the search space regardless of whether the resulting policy is actually better. This was identified after every independently-tuned condition converged on the largest available batch size; switching to a percentile-based objective removes the dependence on pool size.
Every trial for a given condition starts from the same fixed, already-fine-tuned TL checkpoint, so that only the RL steps themselves are repeated across trials.
The optimisation objective is deliberately restricted to the RL reward itself: neither the SurfKit nor the ZINC holdout is used to guide hyperparameter selection, since doing so would constitute a form of validation leakage and invalidate any subsequent rediscovery-rate comparison.
Holdout-based metrics (Section~\ref{sec:methods:evaluation}) are computed only afterward, using the best hyperparameters found.

\subsection{Experimental design: combinatorial evaluation}
\label{sec:methods:design}

Three scoring-function choices are varied independently: (i) whether the structural plausibility term is included (ZincPlausibility on/off); (ii) the uncertainty-handling mode; and (iii) the Pareto-based term added alongside pCMC/SurfTen (none, ParetoBoost, or ParetoGradient).
An initial $2\times4\times3=24$-combination sweep tested all four uncertainty modes (none, Score Modulation, Loss Modulation, and both jointly); Score Modulation and Score+Loss Modulation together did not outperform Loss Modulation alone or no uncertainty handling on renormalised score, and did not add a qualitatively distinct pattern beyond what Loss Modulation already showed, so subsequent combinations retain only the \emph{none} and \emph{Loss Modulation} uncertainty modes, giving $2\times2\times3=12$ combinations.
For a given combination, every scoring component with non-zero weight in Eq.~\ref{eq:sm} is weighted equally, $w_i=1/n_\mathrm{active}$; when Loss Modulation is enabled, the underlying uncertainty components are still computed (so that Eq.~\ref{eq:lm} can read their scores) but are assigned weight zero in the geometric mean.
For each combination, the hyperparameter search of Section~\ref{sec:methods:hpo} is run for 15 trials, followed by 20 independent production replicates (distinct random seeds, step count set by the same oracle-budget rule as the HPO trials) using the best hyperparameters found, all initialised from the same shared TL checkpoint.

\subsection{Evaluation metrics}
\label{sec:methods:evaluation}

All metrics are computed on the pooled set of unique, RDKit-canonicalised SMILES across the 20 production replicates of a given combination (and, where noted, per replicate to obtain a mean and standard deviation), and additionally tracked per RL step to characterise training-time trends rather than only endpoint values.
Validity, uniqueness, novelty, and internal diversity follow standard definitions from the generative-molecular-design benchmarking literature\cite{polykovskiy_molecular_2020}.

\begin{itemize}
    \item \textbf{Validity}: the fraction of generated SMILES strings that RDKit can parse.
    \item \textbf{Uniqueness} and \textbf{novelty}: the fraction of valid molecules that are canonically distinct, and the fraction of unique molecules absent from the transfer-learning trainval set, respectively.
    \item \textbf{Internal diversity}: $1$ minus the mean pairwise Tanimoto similarity (Morgan fingerprints\cite{rogers_extended-connectivity_2010}, radius 2, 2048 bits) over a 2000-molecule subsample of the unique, valid set.
    \item \textbf{Renormalised score}: $\sqrt{s_\mathrm{pCMC}(x)\cdot s_\mathrm{SurfTen}(x)}$ (Eq.~\ref{eq:composite}), computed from the pCMC/SurfTen component scores that REINVENT4 reports for every molecule regardless of which other terms a given combination's objective includes.
    Because pCMC and SurfTen are scored identically across all combinations, this provides a common basis for comparison that is isolated from the effect of the plausibility, uncertainty, and Pareto terms under study.
    \item \textbf{Nearest-neighbour Tanimoto distance to the holdout set}: for each generated molecule, $1$ minus the Tanimoto similarity (Morgan fingerprints) to its single closest match in the SurfKit \emph{holdout} set (Section~\ref{sec:methods:holdout}), averaged over all generated molecules -- distinct from novelty/uniqueness above, which are assessed against the trainval set.
    \item \textbf{SurfKit and ZINC top-$N$ rediscovery rate}: the fraction of the respective holdout (105 molecules for SurfKit, 100 for ZINC; Section~\ref{sec:methods:holdout}) exactly recovered, by canonical SMILES match, among the unique generated molecules.
\end{itemize}

% =====================================================================
% RESULTS
% =====================================================================
\section{Results}
\label{sec:results}

% TODO: populate once the 12-combination production sweep has completed.

% =====================================================================
% DISCUSSION
% =====================================================================
\section{Discussion}
\label{sec:discussion}

% TODO.

% =====================================================================
% CONCLUSION
% =====================================================================
\section{Conclusion}
\label{sec:conclusion}

% TODO.

\section*{Acknowledgment}
% TODO: funding sources, compute resources.

\section*{Data and Software Availability}
All code is available on GitHub \url{https://github.com/hannesohman/REINVENT4Surfactants}.

% =====================================================================
% REFERENCES
% =====================================================================

\bibliographystyle{unsrtnat}
\bibliography{references}

\end{document}
